\documentclass[11pt]{article}
\usepackage[a4paper,margin=0.85in]{geometry}
\usepackage{amsmath,amssymb,booktabs,array}
\usepackage[colorlinks=true,urlcolor=blue,linkcolor=blue]{hyperref}
\setlength{\parindent}{0pt}
\setlength{\parskip}{0.5em}
\setlength{\emergencystretch}{2em}
\newcommand{\dd}{\,\mathrm{d}}
\begin{document}
\begin{center}
{\Large When aggregate flow reveals lane contrast}\\[0.4em]
Ramp-metering review clarification --- 5 October 2026
\end{center}

\section*{1. What is the problem?}
The review asks whether traffic information omitted from the \emph{complete
declared aggregate observation history} can change a useful common meter-rate
decision. A necessary condition is that the compared lane states genuinely look
the same to that observer. Equal road-total mass alone does not meet that condition:
total flow, independently observed speed, counts and past observations also matter.

The slide's ``effective lane factor'' is more precisely the known
\textbf{free-flow mobility coefficient}, not a multiplier on vehicle mass.
For a common cell length \(L\), lane \(\ell\) has
\[
 c_\ell=v_\ell/L,\qquad
 D_\ell(n)=\min(c_\ell n,C_\ell),\qquad
 q_\ell=\min(D_\ell(n_\ell^{\mathrm{after}}),S_{\mathrm{down},\ell}).
\]
Here \(v_\ell\) is the lane free-flow speed, \(c_\ell\) has units inverse time,
\(C_\ell\) is flow capacity, \(D\) is a sending offer, \(S\) is a receiving
limit and \(q\) is actual longitudinal flow rate. ``After'' means after paired
lateral exchange, before longitudinal transmission. These coefficients are
common known parameters across the compared scenarios, not private hidden forcing.
They are also different from the VSL choke-work multiplier or a merge priority.

On an ordinary uncapped interface with ample downstream receiving,
\[
 N=n_1^{\mathrm{after}}+n_2^{\mathrm{after}},\qquad
 Q=q_1+q_2=c_1n_1^{\mathrm{after}}+c_2n_2^{\mathrm{after}}.
\]
\(N\) counts all mass equally; \(Q\) weights the lane masses by their mobility.
Conserving the first sum does not conserve or equalize the second.

\section*{2. A concrete example from the slides}
Take normalized \(c_1=1\), \(c_2=1.5\) per time unit. Both sending offers
are below capacity, downstream receiving is ample and the displayed masses
are post-lateral states of one cell:
\begin{center}
\begin{tabular}{@{}lrrrrrr@{}}
\toprule
State & Lane 1 mass & Lane 2 mass & Total \(N\) & \(q_1\) & \(q_2\) & Total \(Q\)\\
\midrule
A & 2 & 8 & 10 & 2 & 12 & 14\\
B & 6 & 4 & 10 & 6 & 6 & 12\\
\bottomrule
\end{tabular}
\end{center}
Both have ten mass units. State A places more in the more mobile lane, so more
mass can leave this cell per unit time. A road-total flow observation can
distinguish 14 from 12 without reporting separate lane counts.

This is a normalized hand illustration of an observation map, not a calibrated
traffic result, a reachable whole-history pair, a meter intervention or an
optimal-action comparison. Its flow is useful as decision information only if
it was measured at the admitted location/stage and is already available then.

\newpage
\section*{3. Why this matters, and how much can be distinguished}
\textbf{For two lanes the result can be stronger than partial distinction.}
If exact co-supported \(N,Q\) are available, the active branch is known and
\(c_1\ne c_2\), then
\[
 \begin{bmatrix}1&1\\c_1&c_2\end{bmatrix}
 \begin{bmatrix}n_1^{\mathrm{after}}\\n_2^{\mathrm{after}}\end{bmatrix}
 =\begin{bmatrix}N\\Q\end{bmatrix},\qquad
 \det=c_2-c_1\ne0.
\]
Thus the entire two-lane \emph{post-exchange} split is recovered:
\[
 n_1^{\mathrm{after}}=\frac{Q-c_2N}{c_1-c_2},\qquad
 n_2^{\mathrm{after}}=\frac{c_1N-Q}{c_1-c_2}.
\]
For the example, \(N=10,Q=14\) recovers \((2,8)\), and \(N=10,Q=12\)
recovers \((6,4)\). If two free-flow splits differ by \((\delta,-\delta)\),
their total-flow difference is \((c_1-c_2)\delta\): with unequal mobilities,
matching both \(N\) and \(Q\) forces that contrast to vanish.

\textbf{This is a problem for a proposed hidden-information witness, not a
failure of CTM or of conservation.} Such a pair cannot be presented as identical
aggregate observations if this flow sample is among the admitted inputs.
Information can be implicit in a weighted total even when lane labels are absent.
An observer-aware aggregate method can potentially infer that contrast from
public parameters and observations.

\textbf{Information loss and model reduction are different questions.}
A model storing only total mass may fail to predict flow with one constant
aggregate FD, because flow depends on allocation. But a method with access to
total mass \emph{and} flow may know that allocation indirectly. Therefore failure
of a particular mass-only simulator does not prove that the declared strong
aggregate observer lacks the information, or that all aggregate-based controllers
must fail. Reconstruction also does not automatically give a correct traffic
model, optimal meter action or useful gradient.

\textbf{Mass-conserving lane exchange itself illustrates the distinction.}
On the same unimpeded branch, transferring \(h\lambda\) mass from donor \(d\)
to recipient \(t\) preserves the cell total but changes its forward offer by
\[
 Q_{\mathrm{after}}-Q_{\mathrm{before}}
   =h(c_t-c_d)\lambda.
\]
This is the current recomputation rule's local consequence. It is different
from the superseded automatic one-for-one lateral-flow debit, and does not
prove a complete-time benefit or create additional lane capacity.

For more than two lanes, total mass plus total flow generally leaves contrasts
undetermined at a snapshot. Additional speed moments or history can supply more
constraints. This observation does not change the current two-lane task.

\newpage
\section*{4. Which cases resolve the split, and where does the proof stop?}
\renewcommand{\arraystretch}{1.22}
\begin{tabular}{@{}>{\raggedright\arraybackslash}p{0.32\textwidth}>{\raggedright\arraybackslash}p{0.62\textwidth}@{}}
\toprule
Case & Consequence for the information question\\
\midrule
Two lanes; unequal known mobilities; uncapped offers; ample receiving; exact aligned \(N,Q\)
& Unique post-lateral split. This is established under the declared conditions.\\
Equal mobilities on that branch
& \(Q=cN\): flow adds no split information. Speed, other conditions and history still need checking. Equal speeds cannot be imposed just to force a positive result.\\
Capacity, downstream supply or node allocation clips flows
& Accepted flow can conceal underlying sending demand. The linear inverse above no longer applies. This does not prove that speed/history also conceal it.\\
Different measurement locations, averaging windows or stages
& The two measurements may not describe the same state. The inverse cannot be applied without a consistent observation/propagation map.\\
Known earlier state and all accepted per-lane fluxes
& Conservation may reconstruct the current state from completed past observations. This is conditional on complete ramp and boundary balances.\\
Noisy data, uncertain parameters or nearly equal mobilities
& Exact algebra may be practically fragile. That does not establish exact hidden histories under the current idealized review contract.\\
\bottomrule
\end{tabular}

\textbf{A capped-flow illustration.} Keep \(c_1=1,c_2=1.5\), but let
\(C_1=C_2=2\), with enough downstream supply. Splits \((3,7)\) and
\((6,4)\) both have \(N=10\), and both produce \(q_1=q_2=2\), hence
\(Q=4\). All offers are strictly capped. This shows only snapshot ambiguity
in mass/flow: no full speed/history equality or useful action distinction follows.

\textbf{The stage and timing boundary is essential.}
Post-lateral flow identifies a post-lateral state, not necessarily the preceding
start-of-step state. The accepted lateral map can send \((6,4)\) and \((9,1)\)
to the same \((5.4,4.6)\) under the existing rational example. If remaining
state/forcing agree and the control acts only after that shared map, those
erased differences cannot change its action response. Their earlier histories
have not been matched. Unrecoverable information need not be useful information.

Flow caused by a proposed future meter action cannot be supplied to its initial
observer. Reconstruction from a completed previous step needs that step's
identified post-lateral state and actual incoming/outgoing fluxes, including
ramp entry. Current exact ramp/measurement support remains unfinished.

\newpage
\section*{5. How this arose, the evidence, and the remaining gap}
\textbf{Why the audit raised it.} The review retained strong aggregate history,
public lane parameters and complete time accounting. The source has lane-specific
FDs; the chosen task did not assume equal lane mobilities. Auditing which
aggregate channels distinguish a lane split therefore became necessary.
After UG3.B removed the old residual-rate coupling, the observation algebra had
to be rederived for ordinary forward flow from post-exchange occupancy. This
is why matching mass alone was flagged in the slides.

\textbf{Literature support, with its boundary.}
\begin{itemize}
\item \href{https://repository.tudelft.nl/file/File_2c4359e7-5e5a-4651-964f-6f3610ab89ae}{Nagalur Subraveti, Knoop and van Arem, first-order multilane model}:
lane-wise CTM conservation and lateral-transfer structure are relevant
antecedents. Tables 1--2 on printed p.\ 1769 report different calibrated lane
free-flow speeds; for A13 these are 104.6, 98.5 and 87 km/h. Heterogeneity is
therefore a sourced modeling possibility, not a requirement at every site.
\item The same paper, printed pp.\ 1768--1769, distinguishes detector time-mean
speed from model space-mean speed and derives density from detector data.
It does not establish our exact independent cell-mass/passage-speed observer.
\item \href{https://horowitz.me.berkeley.edu/Publications_files/All_papers_numbered/Wright_node_models_transportation_networks_Transportation_research_B_2017.pdf}{Wright et al., node models, Section 3.1}:
demand, receiving supply and priorities determine accepted merge flows.
This grounds the allocation constraints, not lane-state observability or a
congested detector-speed map.
\end{itemize}
The inverse, stage diagnostics and action-erasure argument are \textbf{reviewer
derivations for the adapted RM9 kernel}, not theorems established by those
papers. The sources do not prove that our ramp task is identifiable everywhere.

\textbf{What is already established.} Current claims RM9-C25/C27 provide
conditional post-lateral free-flow reconstruction and prior-step sufficiency;
C26 prevents equating lost pre-state information with useful hidden state.
The strict-branch sensitivity/accounting path is explicit. No full source
algorithm reproduction or numerical traffic experiment supplies these findings.

\textbf{What remains unresolved.}
\begin{enumerate}
\item An adopted observation map: speed meaning, cell versus detector support,
time windows, update stages and available initial/past information.
\item Complete consistent ramp geometry, propagation and mass balances so the
observed history can be related to the state relevant to the next action.
\item Identifiability in capacity/receiving-limited branches under \emph{all}
retained observations, not only a mass/flow snapshot.
\item A feasible pair with identical complete predecision histories but different
useful common meter actions and consequential gradient information, or a
properly scoped aggregate-sufficiency/negative result.
\item Physical observation validity, numerical gradient correctness and independent
gradient usefulness if a later research stage is opened.
\end{enumerate}

\textbf{Current disposition.} UG4 remains open. Its idealized-observer route A
and physical-measurement-first route B remain unapproved. This is an explanation
of existing findings with hand illustrations; it opens no new review track,
observer, parameter choice, simulation or implementation.

\small Local basis: current RM9 separate-resource/observer audit,
claim ledger and source-to-claim audit.
Editable source is retained beside this PDF.
\end{document}
